\section{Relation to action, spectral, and statistical inverse problems}
\label{sec:v5-comparison}
The analysis starts from a discrete variational problem but ends with a
normalized physical probability. Bolotin and Treschev~\cite{Hill} develop
Hill-type identities connecting action Hessians and monodromy in discrete
and continuous Lagrangian systems. Such determinant identities explain
the linear structure used here. The further step in
Theorem~\ref{thm:v4-factorization} is a relative nonlinear determinant
limit with two separated boundary perturbations, followed by full-phase
residual-time integration. Propositions~\ref{prop:v9-cofactor}--\ref{prop:v9-morse-transport}
separate the exact finite identity, trace-class transport and
normalized integration. The additional physical consequence is
Theorems~\ref{thm:v9-compatibility} and~\ref{thm:v9-profile-stability}:
the two same-type scalar laws determine the complete energy profiles
stably in a specified norm, and predict the mixed-type law.
The finite-record compatibility defect in
Corollary~\ref{cor:v9-finite-records} is controlled by bridge length.
It does not follow from just the value of the linear multiplier.

Zelditch~\cite{Zelditch} relates successive localized wave-trace
invariants to analytic boundary geometry. His construction uses
multiple oscillatory boundary integrals and a stationary-phase expansion;
the inverse spectral conclusions concern the specified symmetric classes
of simply connected analytic plane domains, including a symmetry
reversing a bouncing-ball orbit. This is an important predecessor for
the coefficient-to-boundary strategy, not the same observation theorem.
Our coefficients arise from a positive full-phase residual-time integral,
not a Dirichlet or Neumann wave trace. The highest-jet diagonal in
\eqref{eq:v6-finite-diagonal} separates the stationary-action contribution
from the relative interior determinant contribution. Exact analytic germ
recovery and finite-jet coefficient conditioning are kept distinct from
noise-stable global analytic continuation.

Marked-length observations provide a different route to geometry.
B\'alint, De Simoi, Kaloshin and Leguil~\cite{BDKL} recover period-two
curvatures and periodic Lyapunov exponents from the marked length
spectrum of open dispersing billiards. Osterman~\cite{Osterman} studies
analytic rigidity by homoclinic asymptotics and normal forms, including
recovery with two scatterers supplied. De Simoi, Kaloshin and
Leguil~\cite{DSKL} prove marked-length determination of analytic chaotic
billiards with axial symmetries, under symmetry and genericity hypotheses
in a no-eclipse open configuration. Their theorem concerns the geometry
of the table, not merely a specified contact jet.

The data in Corollary~\ref{cor:v5-analytic-rigidity} are instead a selected
onset probability and its offset coefficients under normalized phase
volume. The identical-even contact hypothesis makes its inverse
triangular. Analytic continuation then determines each participating
boundary in its contact frame, but supplies neither a lattice nor the
other obstacles. There is no implication here between equality of these
statistical data and equality of marked length spectra. In particular,
the corollary is not a replacement for the global table determination
in~\cite{DSKL}. Its coefficient formula and its physically realized
finite-jet coordinates identify the precise inverse observation being
used. Proposition~\ref{prop:v5-one-flight} also shows explicitly that
this information is not exclusive to long records.

For periodic finite-horizon Sinai billiards, Finamore and
Leguil~\cite{FinamoreLeguil} introduce an enriched marked length spectrum
and prove isometry from equality of that datum, using geodesic
approximations and CAT(0) geometry. Their enriched observation and
finite-horizon assumptions should not be identified with the selected
count probabilities considered here. Conversely our general boundary-law
construction does not require finite horizon.

The exact exponential-sum step in Theorem~\ref{thm:g-inverse} belongs
to the theory of Prony reconstruction. Batenkov and Yomdin~\cite{BYProny}
analyze local accuracy for confluent Prony systems. The quantitative
result here uses a specific feature of the hyperbolic-sine observation:
its first four functions have the nonzero Wronskian
\eqref{eq:v5-wronskian}. They give regular symmetric coefficient
coordinates even when all three radii coincide. The classical
coefficient-to-root estimate then gives the pairwise exponent $1/3$;
the realized physical path proves its sharpness. No new general root
perturbation theorem or universally improved Prony algorithm is needed.

The source and onset series use a different physical window at each
coefficient. Accordingly their first poles and convergence domains are
consequences of the relative boundary law, not assertions about a
pressure function for all long itineraries. Return-time and compound
Poisson laws such as~\cite{Marklof,CNZ}, and the anisotropic statistical
frameworks in~\cite{DMN,DemersZhang}, address other observations and
asymptotic regimes. The fixed-window companion~\cite{TC} retains its
separate transversality and moving-level argument. The historical
Fourier-inversion analysis~\cite{A2R33} remains available with its own
explicit hypotheses; those hypotheses are not substituted for the
finite-threshold proofs in this article.

Finally, polynomial extrapolation and binomial concentration are
classical ingredients of the acquisition statements. Their uniform use
here depends on the differentiated physical onset expansion. The
self-calibration theorem adds a concrete statistical use of its
quadratic vanishing: square-root probabilities have a simple onset
zero, independent of the unknown normalizer. Its proof accounts for
both the calibration preparations and the subsequent conditional
records, and preserves a positive sampling window on every pilot outcome.
The sufficient cost is not represented as a minimax rate.

The relative boundary law has an experiment-level consequence beyond
coefficient identification. Theorem~\ref{thm:v6-transfer} controls total
variation for the endpoints and residual time, uniformly down to zero
offset, and retains the rare-event mass when unsuccessful preparations
are restored. It permits bounded-risk replacement for growing independent
samples whenever their expected number of successes times $\tau^j$
tends to zero. Its proof uses the quadratic vanishing of the action
error together with the normalized determinant estimate; an absolute
approximation to a vanishing flux would not give this conclusion.

The finite-flight inverse of Theorem~\ref{thm:v6-finite-jets} and the
sensitivity formula \eqref{eq:v6-sensitivity} concern fixed jet order
and varying collision length. The area constraint
\eqref{eq:v6-area-constraint} explains why three leading amplitudes
suffice in the constrained family while four allow an independent
normalizer. Section~\ref{sec:v8-fixed} uses a different fourth datum:
a second first-flight window separates an unknown gap from the three
symmetric shape coordinates in the exact physical model.

The minimax statement in Theorem~\ref{thm:v8-joint-minimax} concerns a
specified binary experiment, not every observation available from a
billiard. Its curvature and timing obstructions are realized by two
different physical pairs. The smooth envelope in
Theorem~\ref{thm:v8-nuisance} removes the supplied exact finite-offset
formula while preserving a positive, fully charged reconstruction
procedure. The earlier supplied-bracket and fixed-bracket procedures
remain valid with their complete proofs. They are neither counterexamples
to the fixed-window benchmark nor general lower bounds for all data.

Finally, the support calculation in Theorem~\ref{thm:v7-exact-tv} is
an elementary disk--ellipse comparison after a common whitening.
Its role is to identify the constant in a physical testing profile,
not to replace the nonlinear relative law. The projection argument in
Theorem~\ref{thm:v8-supercritical} completes that profile when the full
sample's tangent approximation error no longer tends to zero. At fixed
positive offset, the genuinely nonlinear comparison is still supplied
by Theorems~\ref{thm:v8-main-relative} and~\ref{thm:v6-transfer}.

The smooth-envelope optimum in Theorem~\ref{thm:v9-envelope-minimax}
uses the additional strict-slack condition and allows free nuisance
functions. It is not a lower bound obtained from a pair of exact
nonlinear billiard laws. Its attribution to the referee memorandum is
separate from the boundary compatibility and profile inverse developed
here. Neither the count theorem nor its envelope comparison is used
as a substitute for the long-bridge consequence.
